\documentclass{article}
\usepackage[T1,T2A]{fontenc}
\usepackage[utf8]{inputenc}
\usepackage[english,russian]{babel}
\usepackage{xcolor}
\usepackage{amsmath}
\usepackage{amsfonts,bm}
\usepackage{tikz}
\usetikzlibrary{patterns}
\usetikzlibrary{quotes,angles}
\usepackage{relsize}

\title{Решение Задачи №25}
\author{Серёжи Фуканчика\\34 Ю класс}
\date{Декабрь 2019}

\begin{document}

\maketitle

\section{Задача}
В прямоугольном треугольнике $PQR$ из вершины прямого угла $R$ опущена высота $RH$.

Радиусы окружностей, вписанных в треугольники $RHQ$ и $PHR$ равны $6$ и $8$ соответственно. Найдите площадь треугольника.

\section{Решение}
Сделаю рисунок. Нужно нарисовать прямоугольный треугольник и разметить вешины так, чтобы прямой угол был в вершине $R$.

\begin{center}\begin{tikzpicture}
  \draw (0,0) node[left] {$R$} -- (0,2.4)  node[above] {$P$};
  \draw (0,0) -- (2,0) node[above] {$Q$};
  \draw (0,2.4) -- (2,0);

  \draw (0,0) -- (1.1,1.1) node[above right] {$H$};

  \draw (0,0.2) -- (0.2,0.2);
  \draw (0.2,0.2) -- (0.2,0);
\end{tikzpicture}\end{center}

\section{Ответ}

\section{Проверка}

\end{document}
